% I. Анализ на предметната област и литературен преглед.
% Това е разделът, който носи най-голяма тежест в този проект: темата е от
% управленската наука, а не от ежедневната инженерна практика на автора, затова
% всяко твърдение тук стъпва на цитиран източник.
\section{Анализ на предметната област и литературен преглед}
\label{sec:literature}

\subsection{Информационни технологии и подпомагане на управленските решения}

Разделението на управленските решения на структурирани, полуструктурирани и
неструктурирани е въведено от Gorry и Scott Morton~\cite{gorry1971framework} и
остава полезно и днес, защото определя каква роля може да поеме една
информационна система. Структурираното решение се автоматизира изцяло.
Неструктурираното остава изцяло на човека. Полуструктурираното, а изборът на
портфейл от проекти е точно такова, се поддава на изчислителна поддръжка само
частично: изчислима е следствената част, а преценката кои критерии колко тежат
остава управленска. Power~\cite{power2002dss} систематизира класовете системи за
подпомагане на решенията и подчертава същото: системата предоставя модел и
данни, но не поема отговорността за избора.

Оттук следва проектното изискване, което определя целия облик на разработваната
система. Тя не произвежда „правилния отговор“, а прави препоръката проследима:
за всяка избрана позиция трябва да е видимо кое ограничение я е допуснало и коя
оценка я е класирала. \TODO{допълни с втори източник за изискването за
проследимост на препоръката в системите за подпомагане на решенията}

\subsection{Управление и моделиране на бизнес процеси}

Съвременното управление на бизнес процеси се разглежда като жизнен цикъл от
идентификация, откриване, анализ, преработване, внедряване и наблюдение
\cite{dumas2018bpm}. За настоящия проект е съществена фазата на анализ, при
която процесът е вече описан формално и подлежи на количествена оценка по
време, стойност и ресурс. Нотацията BPMN 2.0~\cite{omg2013bpmn} е приетият
стандарт за такова описание и дава семантика на потока от дейности, на
разклоненията и на събитията.

Мрежовият модел на дейностите с продължителности и зависимости, върху който
стъпва изчисляването на най-ранните и най-късните срокове, е по-стар от BPMN и е
известен като метод на критичния път. Той дава на системата две неща наведнъж:
най-ранния възможен срок за завършване на един проект и запаса от време на всяка
негова дейност. Запасът е важен, защото дейност без запас превръща всяко
закъснение в закъснение на целия проект, а това е вход към оценката на риска.
\TODO{цитирай първоизточника за метода на критичния път, Kelley и Walker, 1959,
след като е проверен библиографски}

\subsection{Многокритериален анализ на решенията}

Проектите кандидати не се сравняват по един показател. Стратегическата стойност,
техническият риск, регулаторната необходимост и очакваната възвръщаемост са
несъизмерими и част от тях се измерват само по субективна скала. Методът на
аналитичната йерархия~\cite{saaty1980ahp, saaty1990decision} решава това чрез
попарни сравнения и извежда теглата на критериите от собствен вектор на
матрицата на сравненията, като предлага и мярка за съгласуваност на преценките.
Методът TOPSIS~\cite{hwang1981topsis} подхожда различно: класира алтернативите по
относителна близост до идеалното решение и отдалеченост от най-лошото, при
предварително зададени тегла. Triantaphyllou~\cite{triantaphyllou2000mcdm}
сравнява семейството методи и показва, че различните методи върху едни и същи
данни невинаги дават едно и също класиране. Именно това несъвпадение е причината
двата метода да бъдат реализирани и съпоставени в Раздел~\ref{sec:alternatives},
вместо единият да бъде избран предварително.

\subsection{Избор на портфейл и управление на инвестиционния риск}

Идеята, че портфейлът се оценява като цяло, а не като сбор от независимо добри
позиции, произхожда от Markowitz~\cite{markowitz1952portfolio}. Пренесена върху
портфейл от ИТ проекти, тя означава, че сумарната полза не е единственият
критерий: концентрацията на риск в няколко зависими проекта е свойство на
съвкупността. Стандартът за управление на портфейли~\cite{pmi2017portfolio}
описва управленската рамка около това решение, включително балансирането на
портфейла спрямо стратегическите цели.

От изчислителна страна задачата за избор на подмножество проекти с максимална
сумарна полза при ограничен бюджет е разновидност на задачата за раницата.
Класификацията на вариантите и точните алгоритми са изложени
в~\cite{martello1990knapsack} и~\cite{kellerer2004knapsack}. Задачата е
NP-трудна в общия си вид, но практическите размери в разглеждания случай са
малки, десетки проекти, а не милиони, което прави точното решаване напълно
приложимо. Добавянето на ограничение по капацитет за няколко периода и на
отношения на предхождане между проектите превръща формулировката в
многомерна задача за раницата с ограничения за предхождане.
\TODO{провери и цитирай източник за многомерния вариант с ограничения за
предхождане; не се позовавай на~\cite{kellerer2004knapsack} за твърдение, което
не е проверено в текста му}

\subsection{Изводи от прегледа}

Прегледът дава три опорни точки за проектирането. Първо, системата поддържа
решение, което остава на човека, затова изходът ѝ трябва да е обяснимо, а не
само оптимално множество. Второ, оценяването и подборът са две отделни задачи и
могат да се сменят независимо, което е аргумент за ясна граница между двата
слоя. Трето, изборът на метод за многокритериална оценка не е безразличен за
резултата, затова сравнението между два метода е част от проекта, а не
подробност.
